\documentclass[
    % -- opções da classe memoir --
    12pt,				% tamanho da fonte
    a4paper,            % tamanho do papel
    oneside,			% para impressão em recto e verso. Oposto a twoside
    openright,			% capítulos começam em pág ímpar (insere página vazia caso preciso)
    ruledheader,
    %noindentfirst,
    anapcustomindent,
    sumario = tradicional,
    abntfigtabnum,
    tocpage=plain,
    english,			% idioma adicional para hifenização
	%french,				% idioma adicional para hifenização
	%spanish,			% idioma adicional para hifenização
	brazil,				% o último idioma é o principal do documento
    ]{abntex-ifsuldeminas/abntex2-orientadora} % modificação no abnTeX2 para usar orientadora feminina, dentre outras alterações. Se for orientador realizar a alteração no arquivo.


% ---
% PACOTES
% ---
\usepackage{abntex-ifsuldeminas/abntex-ifsuldeminas}
%Estilo de formatação de capítulos, padrão IFSULDEMINAS

\makeatletter
\newcommand{\thechapterwords}
{ \ifcase \thechapter\or 1\or 2\or 3\or 4\or 5\or 6\or 7\or 8\or 9\or 10\or 11\fi} % se precisar de mais capítulos adicionar conforme a ordem lógica estabelecida.

\def\@makechapterhead#1{%
\vspace*{6cm} % insere o espaço de 6cm da margem, conforme estabelecido.
{\parindent \z@  \reset@font

\bfseries \thechapterwords.  #1\par\nobreak % AQUI ESTÁ O NOME DO CAPÍTULO
\par
\vspace*{10\p@}% pequeno espaçamento após o nome do capítulo
}}

\setsecheadstyle{}
\setsubsecheadstyle{}
\setsubsubsecheadstyle{}
\setparaheadstyle{}
\setsubparaheadstyle{} %modelos dos capítulos

% ---
% Pacotes fundamentais 
% ---
\usepackage{amsmath}
\usepackage{mathspec}
\setmainfont[Ligatures=TeX]{Times New Roman} %Pode ser Arial em ambos
\setsansfont[Ligatures=TeX]{Times New Roman}
\setmonofont[Ligatures=TeX]{Hack}
\setmathrm{Times New Roman}
\setmathfont(Digits,Latin){Times New Roman} % Fonte matemática Times

\usepackage{indentfirst}		            % Identa o primeiro parágrafo de cada seção.
%\usepackage[pdftex]{color,graphicx}			% Inclusão de gráficos e cores
\usepackage{microtype} 			            % para melhorias de justificação
\usepackage{fancyhdr}
\usepackage{hyperref}
\usepackage{gensymb}
\usepackage{lipsum}				% para geração de dummy text
% ---
% Configuração de numeração das figuras, tabelas e equações
%---
\usepackage{chngcntr} %insere o número do capítulo antes no número da figura
\counterwithin{figure}{chapter}
\counterwithin{table}{chapter}
\counterwithin{equation}{chapter}


% ---
% Configuração de fontes (texto e matemática)
%---
%\usepackage{unicode-math}

% ---
% Pacotes adicionais, usados no anexo do modelo de folha de identificação
% ---
\usepackage{multicol}
\usepackage{multirow}
% ---
	
% ---
% Pacotes de citações - Formato ABNT
% biblatex-abnt do Daniel
\usepackage[
    style = abnt, % Sistema alfabético
    %style = abnt-numeric, % Sistema numérico
    %style = abnt-ibid, % Notas de referência
    giveninits, % Abrevia os primeiros nomes na bibliografia
    %backref, % Especifica as páginas em que cada entrada foi citada.
    %citecount, % Além das páginas, especifica quantas vezes cada entrada foi citada.
]{biblatex}
\addbibresource{referencias.bib} % Seu arquivo de bibliografia vai aqui.
\usepackage{csquotes}
% ---

% ---
% Informações de dados para CAPA e FOLHA DE ROSTO
% ---
\titulo{SISTEMA GAD\\Gerenciamento de Atividades Diversificadas} % se for longo pode inserir uma quebra de linha com barras duplas => \\
\autor{Vitor \\ Santos}
\local{Naviraí – MS}
\data{2020}
%\orientador{Prof\textsuperscript{a}. Nome do(a) Orientador(a)}
\coorientador{Prof. Laurentino Augusto Dantas} % descomentar se tiver um coorientador no Trabalho.
\instituicao{Instituto Federal de Mato Grosso do Sul}

\tipotrabalho{Trabalho de Conclusão de Curso } %deixar o espaço no final, para inserir o nome corretamente no preâmbulo.

% O preambulo deve conter o tipo do trabalho, o objetivo, 
% o nome da instituição e a área de concentração 
\preambulo{\imprimirtipotrabalho apresentado como pré-requisito de conclusão do curso de Graduação em Engenharia de Agrimensura e Cartográfica no Instituto Federal de Educação, Ciência e Tecnologia do Sul de Minas Gerais – Campus Inconfidentes, para obtenção do título de Bacharel em Engenharia de Agrimensura e Cartográfica.}
% ---

% ---
% Configurações de aparência do PDF final

% alterando o aspecto da cor azul
\definecolor{blue}{RGB}{29,29,226}

% informações do PDF
\makeatletter
\hypersetup{
     	%pagebackref=true,
		pdftitle={\@title}, 
		pdfauthor={\@author},
    	pdfsubject={\imprimirpreambulo},
	    pdfcreator={LaTeX with abnTeX2, Overleaf},
		pdfkeywords={inserir}{de}{acordo}{com}{a necessidade}, 
		colorlinks=false,       		% false: boxed links; true: colored links
    	linkcolor=blue,          	% color of internal links
    	citecolor=blue,        		% color of links to bibliography
    	filecolor=magenta,      		% color of file links
		urlcolor=blue,
		bookmarksdepth=4
}
\makeatother
% --- 

% --- 
% Espaçamentos entre linhas e parágrafos 
% --- 

% O tamanho do parágrafo é dado por:
\setlength{\parindent}{1.3cm}

% Controle do espaçamento entre um parágrafo e outro:
\setlength{\parskip}{0.2cm}  % tente também \onelineskip

% ---
% compila o indice
% ---
\makeindex
% ---

% ----
% Início do documento
% ----
\begin{document} %-----------------------------------------------------------
% Seleciona o idioma do documento (conforme pacotes do babel)
%\selectlanguage{english}
\selectlanguage{brazil}

% Retira espaço extra obsoleto entre as frases.
\frenchspacing 

% ----------------------------------------------------------
% ELEMENTOS PRÉ-TEXTUAIS
% ----------------------------------------------------------
\pretextual

% ---
% Capa
% ---
\imprimircapa
% ---

% ---
% Folha de rosto
% (o * indica que haverá a ficha bibliográfica)
% ---
\imprimirfolhaderosto
% ---

% ---
% Inserir a ficha bibliografica
% ---

% Isto é um exemplo de Ficha Catalográfica, ou ``Dados internacionais de
% catalogação-na-publicação''. Você pode utilizar este modelo como referência. 
% Porém, provavelmente a biblioteca da sua universidade lhe fornecerá um PDF
% com a ficha catalográfica definitiva após a defesa do trabalho. Quando estiver
% com o documento, salve-o como PDF no diretório do seu projeto e substitua todo
% o conteúdo de implementação deste arquivo pelo comando abaixo:
%
% \begin{fichacatalografica}
%     \includepdf{fig_ficha_catalografica.pdf}
% \end{fichacatalografica}
% ---

% ---
% Inserir folha de aprovação
% ---

% Isto é um exemplo de Folha de aprovação, elemento obrigatório da NBR
% 14724/2011 (seção 4.2.1.3). Você pode utilizar este modelo até a aprovação
% do trabalho. Após isso, substitua todo o conteúdo deste arquivo por uma
% imagem da página assinada pela banca com o comando abaixo:
%
% \begin{folhadeaprovacao}
% \includepdf{folhadeaprovacao_final.pdf}
% \end{folhadeaprovacao}
%
\begin{folhadeaprovacao}

    \begin{center}
    {\bfseries\ABNTEXchapterfont\large\imprimirautor}

    \vspace*{\fill}
    \begin{center}
      \ABNTEXchapterfont\bfseries\large\imprimirtitulo
    \end{center}
    \vspace*{\fill}
    
    %\hspace{.45\textwidth}
    %\begin{minipage}{.5\textwidth}
    %    \imprimirpreambulo
    %\end{minipage}%
    %\vspace*{\fill}

   \bfseries\large{Data de aprovação: DIA de MÊS de ANO}

   \end{center}
   
   %\assinatura{\textbf{Orientadora: \imprimirorientador} \\ \textbf{IFSULDEMINAS - Campus Inconfidentes}} 
   %\assinatura{\textbf{Coorientador: \imprimircoorientador} \\ \textbf{IFMS -  CÂMPUS NAVIRAÍ}}
   \assinatura{\textbf{Prof. Laurentino Augusto Dantas} \\ \textbf{IFMS - Naviraí}}
   %\assinatura{\textbf{Professor} \\ Convidado 3}
   %\assinatura{\textbf{Professor} \\ Convidado 4}
    
    \vspace*{1cm}
    
   %\begin{center}
    %\vspace*{0.5cm}
    %{\large\imprimirlocal}
    %\par
    %{\large\imprimirdata}
    %\vspace*{1cm}
   %\end{center}
  
\end{folhadeaprovacao}
% ---

% ---
% Dedicatória
% ---
\begin{dedicatoria}
   \vspace*{\fill}
   \centering
   \noindent
   \textit{ Escreva aqui \\ sua dedicatória \\ se desejar. } \vspace*{\fill}
\end{dedicatoria}
% ---

% ---
% Agradecimentos
% ---
\begin{agradecimentos}

Descrever aqui os agradecimentos.


\end{agradecimentos}
% ---

% ---
% Epígrafe
% ---
\begin{epigrafe}
    \vspace*{\fill}
	\begin{flushright}
		\textit{``Não se amoldem ao padrão deste mundo, \\
		mas transformem-se pela renovação da sua mente, \\
		para que sejam capazes de experimentar e comprovar \\
		a boa, agradável e perfeita vontade de Deus.''\\
		(Bíblia Sagrada. Romanos 12, 2)}
	\end{flushright}
\end{epigrafe}
% ---

% ---
% RESUMOS
% ---

% resumo em português
\setlength{\absparsep}{18pt} % ajusta o espaçamento dos parágrafos do resumo
\begin{resumo}
 Segundo a NBR6028 2003, o resumo deve ressaltar o
 objetivo, o método, os resultados e as conclusões do documento. A ordem e a extensão
 destes itens dependem do tipo de resumo (informativo ou indicativo) e do
 tratamento que cada item recebe no documento original. O resumo deve ser
 precedido da referência do documento, com exceção do resumo inserido no
 próprio documento. (\ldots) As palavras-chave devem figurar logo abaixo do
 resumo, antecedidas da expressão ``Palavras-chave:'', separadas entre si por
 ponto e finalizadas também por ponto.

 \textbf{Palavras-chave}: \LaTeX. Palavras em maiúsculo. Editoração de texto.
\end{resumo}

% resumo em inglês
\begin{resumo}[ABSTRACT]
 \begin{otherlanguage*}{english}
   This is the english abstract.

   \vspace{\onelineskip}
 
   \noindent 
   \textbf{Keywords}: \LaTeX. Capitalized words. Text editoration.
 \end{otherlanguage*}
\end{resumo}

% ---
% inserir lista de figuras
% ---
\pdfbookmark[0]{\listfigurename}{lof}
\listoffigures*
\cleardoublepage
% ---

% ---
% inserir lista de tabelas
% ---
\pdfbookmark[0]{\listtablename}{lot}
\listoftables*
\cleardoublepage
% ---

% ---
% inserir lista de abreviações e siglas
% ---
\begin{siglas}
  \item [SIGLA] Significado por extenso. 
  \item [IBGE] Instituto Brasileiro de Geografia e Estatística
\end{siglas}
% ---

% ---
% inserir lista de símbolos
% ---
\begin{simbolos}
  \item[$ \Gamma $] Letra grega Gama
  \item[$ \Lambda $] Lambda
  \item[$ \zeta $] Letra grega minúscula zeta
  \item[$ \in $] Pertence
\end{simbolos}
% ---

% ---
% inserir o sumario
% ---
\renewcommand{\contentsname}{SUM\'ARIO} % para deixar o nome do sumário em maiúsculo.
\pdfbookmark[0]{\contentsname}{toc}
\tableofcontents*
\cleardoublepage
% ---

% ----------------------------------------------------------
% ELEMENTOS TEXTUAIS
% ----------------------------------------------------------
\textual

% ----------------------------------------------------------
% Introdução (no IFSULDEMINAS, precisa ser numerada como os outros capítulos.)
% ----------------------------------------------------------
\chapter{INTRODUÇÃO}
\thispagestyle{empty} % para não numerar a página da introdução.
    
    \lipsum[1-4] %dummy text. Deverá ser removido na versão final.
    
% ----------------------------------------------------------


% ---
% Capitulo de revisão de literatura
% ---
\chapter{REVISÃO DE LITERATURA}
    
    Exemplo de citação na linha. Deve ser utilizado o comando conforme o código fonte. Segundo \textcite{Andrade2003}, \textbf{Fotogrametria} é \dots
    
    Exemplo de citação ao final de um parágrafo. Também deve ser inserido conforme explicação do código fonte. \cite{Monico2007}
    
    Para maiores informações sobre como fazer as citações, acessar o documento \href{http://mirrors.ctan.org/macros/latex/contrib/biblatex-contrib/biblatex-abnt/doc/biblatex-abnt.pdf}{\color{blue} \textbf{biblatex}}, de Daniel Marques.
    
    \textbf{Lembre-se de remover todos os \textit{dummy text} inseridos neste template.}
    
    Ambiente de equação numerada, de acordo com o capítulo.
    
    \begin{equation}
        a^2 + b^2 = c^2
    \end{equation}
    
    Exemplo de citação direta longa, no formato ABNT e seus devidos recuos:
    
    \begin{citacao}
        \lipsum[1] \cite{CoelhoFilho2007}
    \end{citacao}
    
    \lipsum[1-4]
    
% ----------------------------------------------------------

% ---
% Capitulo de objetivos
% ---
\chapter{OBJETIVOS}
% ---
\section{OBJETIVOS GERAIS}
\lipsum[1]

\section{OBJETIVOS ESPECÍFICOS}
\lipsum[2]

% ----------------------------------------------------------

% ---
% Materiais e métodos
% ---
\chapter{MATERIAIS E MÉTODOS}

    \lipsum[1]
    
    \section{EXEMPLO DE SEÇÃO}
    
        \lipsum[2]
        
        \subsection{EXEMPLO DE SUBSEÇÃO}
        
            \lipsum[3]
            
            \subsubsection{EXEMPLO DE ``SUBSUBSEÇÃO''} %mostra como inserir aspas de maneira correta, abrindo e fechando.
            
                \lipsum[4]

% ---
% Materiais e métodos
% ---
\chapter{RESULTADOS}

\lipsum[27-30]

% ----------------------------------------------------------
% Finaliza a parte no bookmark do PDF
% para que se inicie o bookmark na raiz
% e adiciona espaço de parte no Sumário
% ----------------------------------------------------------
%\phantompart

% ---
% Conclusão
% ---

\chapter{CONCLUSÃO}

\lipsum[31-33]

% ----------------------------------------------------------
% ELEMENTOS PÓS-TEXTUAIS
% ----------------------------------------------------------
\postextual
% ----------------------------------------------------------

% ----------------------------------------------------------
% Referências bibliográficas
% ----------------------------------------------------------

\printbibliography[title={REFERÊNCIAS}] % para inserir o nome das referências em maiúsculo.

% ----------------------------------------------------------
% Apêndices
% ----------------------------------------------------------

% ---
% Inicia os apêndices
% ---
%\begin{apendicesenv}

% Imprime uma página indicando o início dos apêndices
% \partapendices

% ----------------------------------------------------------
%\chapter{EXEMPLO DE APÊNDICE}

%\end{apendicesenv}
% ---


% ----------------------------------------------------------
% Anexos
% ----------------------------------------------------------

% ---
% Inicia os anexos
% ---
%\begin{anexosenv}

% Imprime uma página indicando o início dos anexos
%\partanexos

%\chapter{SE PRECISAR DOS ANEXOS}

%\end{anexosenv}

\end{document}